%%%
%%% Autogenerated file - do not edit 
%%%
\chapter{Heaps}
\chaplabel{heaps}
In this chapter, we discuss two implementations of the extremely useful
priority Queue data structure.  Both of these structures are a special
kind of binary tree called a \emph{heap}, 
\index{heap}%
\index{binary heap}%
\index{heap!binary}%
which means ``a disorganized
pile.''  This is in contrast to binary search trees that can be thought
of as a highly organized pile.
The first heap implementation uses an array to simulate a complete binary
tree.  This very fast implementation is the basis of one of the fastest
known sorting algorithms, namely heapsort (see \secref{heapsort}).
The second implementation is based on more flexible binary trees.
It supports a \ensuremath{\ensuremath{\mathrm{meld}(\ensuremath{\mathit{h}})}} operation that allows the priority queue to
absorb the elements of a second priority queue \ensuremath{\ensuremath{\ensuremath{\mathit{h}}}}.
\section{BinaryHeap: An Implicit Binary Tree}
\seclabel{binaryheap}
\index{BinaryHeap@BinaryHeap}%
Our first implementation of a (priority) Queue is based on a technique
that is over four hundred years old.  \emph{Eytzinger's method}
\index{Eytzinger's method}%
allows us
to represent a complete binary tree as an array by laying out the nodes
of the tree in breadth-first order (see \secref{bintree:traversal}).
In this way, the root is stored at position 0, the root's left child is
stored at position 1, the root's right child at position 2, the left
child of the left child of the root is stored at position 3, and so
on. See \figref{eytzinger}.
\begin{figure}
  \begin{center}
    \includegraphics[scale=0.90909]{figs-python/eytzinger}
  \end{center}
  \caption{Eytzinger's method represents a complete binary tree as an array.}
  \figlabel{eytzinger}
\end{figure}
If we apply Eytzinger's method to a sufficiently large tree, some
patterns emerge.  The left child of the node at index \ensuremath{\ensuremath{\ensuremath{\mathit{i}}}} is at index
$\ensuremath{\ensuremath{\mathrm{left}(\ensuremath{\mathit{i}})}}=2\ensuremath{\ensuremath{\ensuremath{\mathit{i}}}}+1$ and the right child of the node at index \ensuremath{\ensuremath{\ensuremath{\mathit{i}}}} is at
index $\ensuremath{\ensuremath{\mathrm{right}(\ensuremath{\mathit{i}})}}=2\ensuremath{\ensuremath{\ensuremath{\mathit{i}}}}+2$.  The parent of the node at index \ensuremath{\ensuremath{\ensuremath{\mathit{i}}}} is at
index $\ensuremath{\ensuremath{\mathrm{parent}(\ensuremath{\mathit{i}})}}=(\ensuremath{\ensuremath{\ensuremath{\mathit{i}}}}-1)/2$.
%importing \codeimport{ods/BinaryHeap.left(i).right(i).parent(i)}: 
\begin{oframed}
\begin{flushleft}
\hspace*{1em} \ensuremath{\mathrm{left}(\ensuremath{\mathit{i}})}\\
\hspace*{1em} \hspace*{1em} {\color{black} \textbf{return}} \ensuremath{2\cdot \ensuremath{\mathit{i}} + 1}\\
\ \\
\hspace*{1em} \ensuremath{\mathrm{right}(\ensuremath{\mathit{i}})}\\
\hspace*{1em} \hspace*{1em} {\color{black} \textbf{return}} \ensuremath{2\cdot (\ensuremath{\mathit{i}}+1)}\\
\ \\
\hspace*{1em} \ensuremath{\mathrm{parent}(\ensuremath{\mathit{i}})}\\
\hspace*{1em} \hspace*{1em} {\color{black} \textbf{return}} \ensuremath{(\ensuremath{\mathit{i}}-1)\bdiv 2}\\
\end{flushleft}
\end{oframed}
A BinaryHeap uses this technique to implicitly represent a complete
binary tree in which the elements are \emph{heap-ordered}:
\index{heap-ordered binary tree}%
\index{binary tree!heap-ordered}%
\index{heap order}%
The value
stored at any index \ensuremath{\ensuremath{\ensuremath{\mathit{i}}}} is not smaller than the value stored at index
\ensuremath{\ensuremath{\mathrm{parent}(\ensuremath{\mathit{i}})}}, with the exception of the root value, $\ensuremath{\ensuremath{\ensuremath{\mathit{i}}}}=0$.  It follows
that the smallest value in the priority Queue is therefore stored at
position 0 (the root).
In a BinaryHeap, the \ensuremath{\ensuremath{\ensuremath{\mathit{n}}}} elements are stored in an array \ensuremath{\ensuremath{\ensuremath{\mathit{a}}}}:
%importing \codeimport{ods/BinaryHeap.a.n}: 
\begin{oframed}
\begin{flushleft}
\hspace*{1em} \ensuremath{\mathrm{initialize}()}\\
\hspace*{1em} \hspace*{1em} \ensuremath{\ensuremath{\mathit{a}} \gets  \ensuremath{\mathrm{new\_array}(1)}}\\
\hspace*{1em} \hspace*{1em} \ensuremath{\ensuremath{\mathit{n}} \gets  \ensuremath{0}}\\
\end{flushleft}
\end{oframed}
Implementing the \ensuremath{\ensuremath{\mathrm{add}(\ensuremath{\mathit{x}})}} operation is fairly straightforward.  As with
all array-based structures, we first check to see if \ensuremath{\ensuremath{\ensuremath{\mathit{a}}}} is full (by checking if $\ensuremath{\ensuremath{\mathrm{length}(\ensuremath{\mathit{a}})}}=\ensuremath{\ensuremath{\ensuremath{\mathit{n}}}}$) and, if so, we grow \ensuremath{\ensuremath{\ensuremath{\mathit{a}}}}.  Next, we place \ensuremath{\ensuremath{\ensuremath{\mathit{x}}}} at location
\ensuremath{\ensuremath{\ensuremath{\mathit{a}}[\ensuremath{\mathit{n}}]}} and increment \ensuremath{\ensuremath{\ensuremath{\mathit{n}}}}.  At this point, all that remains is to ensure
that we maintain the heap property.  We do this by repeatedly swapping
\ensuremath{\ensuremath{\ensuremath{\mathit{x}}}} with its parent until \ensuremath{\ensuremath{\ensuremath{\mathit{x}}}} is no longer smaller than its parent.
See \figref{heap-insert}.
%importing \codeimport{ods/BinaryHeap.add(x).bubbleUp(i)}: 
\begin{oframed}
\begin{flushleft}
\hspace*{1em} \ensuremath{\mathrm{add}(\ensuremath{\mathit{x}})}\\
\hspace*{1em} \hspace*{1em} {\color{black} \textbf{if}} \ensuremath{\mathrm{length}(\ensuremath{\mathit{a}}) < \ensuremath{\mathit{n}}+1} {\color{black} \textbf{then}} \\
\hspace*{1em} \hspace*{1em} \hspace*{1em} \ensuremath{\mathrm{resize}()}\\
\hspace*{1em} \hspace*{1em} \ensuremath{\ensuremath{\mathit{a}}[\ensuremath{n}] \gets  \ensuremath{x}}\\
\hspace*{1em} \hspace*{1em} \ensuremath{\ensuremath{\mathit{n}} \gets  \ensuremath{\ensuremath{\mathit{n}} + 1}}\\
\hspace*{1em} \hspace*{1em} \ensuremath{\mathrm{bubble\_up}(\ensuremath{\mathit{n}}-1)}\\
\hspace*{1em} \hspace*{1em} {\color{black} \textbf{return}} \ensuremath{\ensuremath{\mathit{true}}}\\
\ \\
\hspace*{1em} \ensuremath{\mathrm{bubble\_up}(\ensuremath{\mathit{i}})}\\
\hspace*{1em} \hspace*{1em} \ensuremath{\ensuremath{\mathit{p}} \gets  \ensuremath{\mathrm{parent}(\ensuremath{\mathit{i}})}}\\
\hspace*{1em} \hspace*{1em} {\color{black} \textbf{while}} \ensuremath{\ensuremath{\mathit{i}} > 0} {\color{black} \textbf{and}} \ensuremath{\ensuremath{\mathit{a}}[\ensuremath{\mathit{i}}] < \ensuremath{\mathit{a}}[\ensuremath{\mathit{p}}]} {\color{black} \textbf{do}} \\
\hspace*{1em} \hspace*{1em} \hspace*{1em} \ensuremath{\ensuremath{\mathit{a}}[\ensuremath{i}], \ensuremath{\ensuremath{\mathit{a}}[\ensuremath{\mathit{p}}] \gets  \ensuremath{\mathit{a}}[\ensuremath{\mathit{p}}], \ensuremath{\mathit{a}}[\ensuremath{\mathit{i}}]}}\\
\hspace*{1em} \hspace*{1em} \hspace*{1em} \ensuremath{\ensuremath{\mathit{i}} \gets  \ensuremath{p}}\\
\hspace*{1em} \hspace*{1em} \hspace*{1em} \ensuremath{\ensuremath{\mathit{p}} \gets  \ensuremath{\mathrm{parent}(\ensuremath{\mathit{i}})}}\\
\end{flushleft}
\end{oframed}
\begin{figure}
  \begin{center}
    \includegraphics[height=\QuarterHeightScaleIfNeeded]{figs-python/heap-insert-1} \\
    \includegraphics[height=\QuarterHeightScaleIfNeeded]{figs-python/heap-insert-2} \\
    \includegraphics[height=\QuarterHeightScaleIfNeeded]{figs-python/heap-insert-3} \\
    \includegraphics[height=\QuarterHeightScaleIfNeeded]{figs-python/heap-insert-4} \\
  \end{center}
  \caption[Adding to a BinaryHeap]{Adding the value 6 to a BinaryHeap.}
  \figlabel{heap-insert}
\end{figure}
Implementing the \ensuremath{\ensuremath{\mathrm{remove}()}} operation, which removes the smallest value
from the heap, is a little trickier.  We know where the smallest value is
(at the root), but we need to replace it after we remove it and ensure
that we maintain the heap property.
The easiest way to do this is to replace the root with the value \ensuremath{\ensuremath{\ensuremath{\mathit{a}}[\ensuremath{\mathit{n}}-1]}}, delete
that value, and decrement \ensuremath{\ensuremath{\ensuremath{\mathit{n}}}}.  Unfortunately, the new root element is now
probably not the smallest element, so it needs to be moved downwards.
We do this by repeatedly comparing this element to its two children.
If it is the smallest of the three then we are done.  Otherwise, we swap
this element with the smallest of its two children and continue.
%importing \codeimport{ods/BinaryHeap.remove().trickleDown(i)}: 
\begin{oframed}
\begin{flushleft}
\hspace*{1em} \ensuremath{\mathrm{remove}()}\\
\hspace*{1em} \hspace*{1em} \ensuremath{\ensuremath{\mathit{x}} \gets  \ensuremath{\ensuremath{\mathit{a}}[0]}}\\
\hspace*{1em} \hspace*{1em} \ensuremath{\ensuremath{\mathit{a}}[\ensuremath{0}] \gets  \ensuremath{\ensuremath{\mathit{a}}[\ensuremath{\mathit{n}}-1]}}\\
\hspace*{1em} \hspace*{1em} \ensuremath{\ensuremath{\mathit{n}} \gets  \ensuremath{\ensuremath{\mathit{n}} - 1}}\\
\hspace*{1em} \hspace*{1em} \ensuremath{\mathrm{trickle\_down}(0)}\\
\hspace*{1em} \hspace*{1em} {\color{black} \textbf{if}} \ensuremath{3\cdot \ensuremath{\mathit{n}} < \mathrm{length}(\ensuremath{\mathit{a}})} {\color{black} \textbf{then}} \\
\hspace*{1em} \hspace*{1em} \hspace*{1em} \ensuremath{\mathrm{resize}()}\\
\hspace*{1em} \hspace*{1em} {\color{black} \textbf{return}} \ensuremath{\ensuremath{\mathit{x}}}\\
\ \\
\hspace*{1em} \ensuremath{\mathrm{trickle\_down}(\ensuremath{\mathit{i}})}\\
\hspace*{1em} \hspace*{1em} {\color{black} \textbf{while}} \ensuremath{\ensuremath{\mathit{i}} \ge 0} {\color{black} \textbf{do}} \\
\hspace*{1em} \hspace*{1em} \hspace*{1em} \ensuremath{\ensuremath{\mathit{j}} \gets  \ensuremath{-1}}\\
\hspace*{1em} \hspace*{1em} \hspace*{1em} \ensuremath{\ensuremath{\mathit{r}} \gets  \ensuremath{\mathrm{right}(\ensuremath{\mathit{i}})}}\\
\hspace*{1em} \hspace*{1em} \hspace*{1em} {\color{black} \textbf{if}} \ensuremath{\ensuremath{\mathit{r}} < n} {\color{black} \textbf{and}} \ensuremath{\ensuremath{\mathit{a}}[\ensuremath{\mathit{r}}] < \ensuremath{\mathit{a}}[\ensuremath{\mathit{i}}]} {\color{black} \textbf{then}} \\
\hspace*{1em} \hspace*{1em} \hspace*{1em} \hspace*{1em} \ensuremath{\ensuremath{\mathit{\ell}} \gets  \ensuremath{\mathrm{left}(\ensuremath{\mathit{i}})}}\\
\hspace*{1em} \hspace*{1em} \hspace*{1em} \hspace*{1em} {\color{black} \textbf{if}} \ensuremath{\ensuremath{\mathit{a}}[\ensuremath{\mathit{\ell}}] < \ensuremath{\mathit{a}}[\ensuremath{\mathit{r}}]} {\color{black} \textbf{then}} \\
\hspace*{1em} \hspace*{1em} \hspace*{1em} \hspace*{1em} \hspace*{1em} \ensuremath{\ensuremath{\mathit{j}} \gets  \ensuremath{\ell}}\\
\hspace*{1em} \hspace*{1em} \hspace*{1em} \hspace*{1em} {\color{black} \textbf{else}} \\
\hspace*{1em} \hspace*{1em} \hspace*{1em} \hspace*{1em} \hspace*{1em} \ensuremath{\ensuremath{\mathit{j}} \gets  \ensuremath{r}}\\
\hspace*{1em} \hspace*{1em} \hspace*{1em} {\color{black} \textbf{else}} \\
\hspace*{1em} \hspace*{1em} \hspace*{1em} \hspace*{1em} \ensuremath{\ensuremath{\mathit{\ell}} \gets  \ensuremath{\mathrm{left}(\ensuremath{\mathit{i}})}}\\
\hspace*{1em} \hspace*{1em} \hspace*{1em} \hspace*{1em} {\color{black} \textbf{if}} \ensuremath{\ensuremath{\mathit{\ell}} < n} {\color{black} \textbf{and}} \ensuremath{\ensuremath{\mathit{a}}[\ensuremath{\mathit{\ell}}] < \ensuremath{\mathit{a}}[\ensuremath{\mathit{i}}]} {\color{black} \textbf{then}} \\
\hspace*{1em} \hspace*{1em} \hspace*{1em} \hspace*{1em} \hspace*{1em} \ensuremath{\ensuremath{\mathit{j}} \gets  \ensuremath{\ell}}\\
\hspace*{1em} \hspace*{1em} \hspace*{1em} {\color{black} \textbf{if}} \ensuremath{\ensuremath{\mathit{j}} \ge 0} {\color{black} \textbf{then}} \\
\hspace*{1em} \hspace*{1em} \hspace*{1em} \hspace*{1em} \ensuremath{\ensuremath{\mathit{a}}[\ensuremath{j}], \ensuremath{\ensuremath{\mathit{a}}[\ensuremath{\mathit{i}}] \gets\hspace*{1em}   \ensuremath{\mathit{a}}[\ensuremath{\mathit{i}}], \ensuremath{\mathit{a}}[\ensuremath{\mathit{j}}]}}\\
\hspace*{1em} \hspace*{1em} \hspace*{1em} \ensuremath{\ensuremath{\mathit{i}} \gets  \ensuremath{j}}\\
\end{flushleft}
\end{oframed}
\begin{figure}
  \begin{center}
    \includegraphics[height=\QuarterHeightScaleIfNeeded]{figs-python/heap-remove-1} \\
    \includegraphics[height=\QuarterHeightScaleIfNeeded]{figs-python/heap-remove-2} \\
    \includegraphics[height=\QuarterHeightScaleIfNeeded]{figs-python/heap-remove-3} \\
    \includegraphics[height=\QuarterHeightScaleIfNeeded]{figs-python/heap-remove-4} \\
  \end{center}
  \caption[Removing from a BinaryHeap]{Removing the minimum value, 4, from a BinaryHeap.}
  \figlabel{heap-remove}
\end{figure}
As with other array-based structures, we will ignore the time spent in
calls to \ensuremath{\ensuremath{\mathrm{resize}()}}, since these can be accounted for using the amortization
argument from \lemref{arraystack-amortized}.  The running times of
both \ensuremath{\ensuremath{\mathrm{add}(\ensuremath{\mathit{x}})}} and \ensuremath{\ensuremath{\mathrm{remove}()}} then depend on the height of the (implicit)
binary tree.  Luckily, this is a \emph{complete}
\index{binary tree!complete}%
\index{complete binary tree}%
binary tree;  every level
except the last has the maximum possible number of nodes.  Therefore,
if the height of this tree is $h$, then it has at least $2^h$ nodes.
Stated another way
\[
  \ensuremath{\ensuremath{\ensuremath{\mathit{n}}}} \ge 2^h \enspace .
\]  
Taking logarithms on both sides of this equation gives
\[
   h \le \log \ensuremath{\ensuremath{\ensuremath{\mathit{n}}}} \enspace .
\]
Therefore, both the \ensuremath{\ensuremath{\mathrm{add}(\ensuremath{\mathit{x}})}} and \ensuremath{\ensuremath{\mathrm{remove}()}} operation run in $O(\log \ensuremath{\ensuremath{\ensuremath{\mathit{n}}}})$ time.
\subsection{Summary}
The following theorem summarizes the performance of a BinaryHeap:
\begin{thm}\thmlabel{binaryheap}
  A BinaryHeap implements the (priority) Queue interface.  Ignoring
  the cost of calls to \ensuremath{\ensuremath{\mathrm{resize}()}}, a BinaryHeap supports the operations
  \ensuremath{\ensuremath{\mathrm{add}(\ensuremath{\mathit{x}})}} and \ensuremath{\ensuremath{\mathrm{remove}()}} in $O(\log \ensuremath{\ensuremath{\ensuremath{\mathit{n}}}})$ time per operation.
  Furthermore, beginning with an empty BinaryHeap, any sequence of $m$
  \ensuremath{\ensuremath{\mathrm{add}(\ensuremath{\mathit{x}})}} and \ensuremath{\ensuremath{\mathrm{remove}()}} operations results in a total of $O(m)$
  time spent during all calls to \ensuremath{\ensuremath{\mathrm{resize}()}}.
\end{thm}
\section{MeldableHeap: A Randomized Meldable Heap}
\seclabel{meldableheap}
\index{MeldableHeap@MeldableHeap}%
In this section, we describe the MeldableHeap, a priority Queue
implementation in which the underlying structure is also a heap-ordered
binary tree.  However, unlike a BinaryHeap in which the underlying
binary tree is completely defined by the number of elements, there
are no restrictions on the shape of the binary tree that underlies
a MeldableHeap; anything goes.
The \ensuremath{\ensuremath{\mathrm{add}(\ensuremath{\mathit{x}})}} and \ensuremath{\ensuremath{\mathrm{remove}()}} operations in a MeldableHeap are
implemented in terms of the \ensuremath{\ensuremath{\mathrm{merge}(\ensuremath{\mathit{h_1}},\ensuremath{\mathit{h_2}})}} operation.  This operation
takes two heap nodes \ensuremath{\ensuremath{\ensuremath{\mathit{h_1}}}} and \ensuremath{\ensuremath{\ensuremath{\mathit{h_2}}}} and merges them, returning a heap
node that is the root of a heap that contains all elements in the subtree
rooted at \ensuremath{\ensuremath{\ensuremath{\mathit{h_1}}}} and all elements in the subtree rooted at \ensuremath{\ensuremath{\ensuremath{\mathit{h_2}}}}.
The nice thing about a \ensuremath{\ensuremath{\mathrm{merge}(\ensuremath{\mathit{h_1}},\ensuremath{\mathit{h_2}})}} operation is that it can be
defined recursively. See \figref{meldable-merge}.  If either \ensuremath{\ensuremath{\ensuremath{\mathit{h_1}}}} or
\ensuremath{\ensuremath{\ensuremath{\mathit{h_2}}}} is \ensuremath{\ensuremath{\ensuremath{\mathit{nil}}}}, then we are merging with an empty set, so we return \ensuremath{\ensuremath{\ensuremath{\mathit{h_2}}}}
or \ensuremath{\ensuremath{\ensuremath{\mathit{h_1}}}}, respectively.  Otherwise, assume $\ensuremath{\ensuremath{\ensuremath{\mathit{h_1}}.\ensuremath{\mathit{x}}}} \le \ensuremath{\ensuremath{\ensuremath{\mathit{h_2}}.\ensuremath{\mathit{x}}}}$ since,
if $\ensuremath{\ensuremath{\ensuremath{\mathit{h_1}}.\ensuremath{\mathit{x}}}} > \ensuremath{\ensuremath{\ensuremath{\mathit{h_2}}.\ensuremath{\mathit{x}}}}$, then we can reverse the roles of \ensuremath{\ensuremath{\ensuremath{\mathit{h_1}}}} and \ensuremath{\ensuremath{\ensuremath{\mathit{h_2}}}}.
Then we know that the root of the merged heap will contain \ensuremath{\ensuremath{\ensuremath{\mathit{h_1}}.\ensuremath{\mathit{x}}}}, and
we can recursively merge \ensuremath{\ensuremath{\ensuremath{\mathit{h_2}}}} with \ensuremath{\ensuremath{\ensuremath{\mathit{h_1}}.\ensuremath{\mathit{left}}}} or \ensuremath{\ensuremath{\ensuremath{\mathit{h_1}}.\ensuremath{\mathit{right}}}}, as we wish.
This is where randomization comes in, and we toss a coin to decide
whether to merge \ensuremath{\ensuremath{\ensuremath{\mathit{h_2}}}} with \ensuremath{\ensuremath{\ensuremath{\mathit{h_1}}.\ensuremath{\mathit{left}}}} or \ensuremath{\ensuremath{\ensuremath{\mathit{h_1}}.\ensuremath{\mathit{right}}}}:
%importing \codeimport{ods/MeldableHeap.merge(h1,h2)}: 
\begin{oframed}
\begin{flushleft}
\hspace*{1em} \ensuremath{\mathrm{merge}(\ensuremath{\mathit{h_1}}, \ensuremath{\mathit{h_2}})}\\
\hspace*{1em} \hspace*{1em} {\color{black} \textbf{if}} \ensuremath{\ensuremath{\mathit{h_1}} \eq nil} {\color{black} \textbf{then}}  {\color{black} \textbf{return}} \ensuremath{\ensuremath{\mathit{h_2}}}\\
\hspace*{1em} \hspace*{1em} {\color{black} \textbf{if}} \ensuremath{\ensuremath{\mathit{h_2}} \eq nil} {\color{black} \textbf{then}}  {\color{black} \textbf{return}} \ensuremath{\ensuremath{\mathit{h_1}}}\\
\hspace*{1em} \hspace*{1em} {\color{black} \textbf{if}} \ensuremath{\ensuremath{\mathit{h_2}}.\ensuremath{\mathit{x}} < \ensuremath{\mathit{h_1}}.x} {\color{black} \textbf{then}}  \ensuremath{(\ensuremath{\ensuremath{\mathit{h_1}}, h2}) \gets  \ensuremath{(\ensuremath{\mathit{h_2}}, \ensuremath{\mathit{h_1}})}}\\
\hspace*{1em} \hspace*{1em} {\color{black} \textbf{if}} \ensuremath{\mathrm{random\_bit}()} {\color{black} \textbf{then}} \\
\hspace*{1em} \hspace*{1em} \hspace*{1em} \ensuremath{\ensuremath{\mathit{h_1}}.\ensuremath{\ensuremath{\mathit{left}} \gets  \mathrm{merge}(\ensuremath{\mathit{h_1}}.\ensuremath{\mathit{left}}, \ensuremath{\mathit{h_2}})}}\\
\hspace*{1em} \hspace*{1em} \hspace*{1em} \ensuremath{\ensuremath{\mathit{h_1}}.\ensuremath{\ensuremath{\mathit{left}}.\ensuremath{\mathit{parent}} \gets  h1}}\\
\hspace*{1em} \hspace*{1em} {\color{black} \textbf{else}} \\
\hspace*{1em} \hspace*{1em} \hspace*{1em} \ensuremath{\ensuremath{\mathit{h_1}}.\ensuremath{\ensuremath{\mathit{right}} \gets  \mathrm{merge}(\ensuremath{\mathit{h_1}}.\ensuremath{\mathit{right}}, \ensuremath{\mathit{h_2}})}}\\
\hspace*{1em} \hspace*{1em} \hspace*{1em} \ensuremath{\ensuremath{\mathit{h_1}}.\ensuremath{\ensuremath{\mathit{right}}.\ensuremath{\mathit{parent}} \gets  h1}}\\
\hspace*{1em} \hspace*{1em} {\color{black} \textbf{return}} \ensuremath{\ensuremath{\mathit{h_1}}}\\
\end{flushleft}
\end{oframed}
\begin{figure}
  \centering{\includegraphics[width=\ScaleIfNeeded]{figs-python/meldable-merge}}
  \caption[Merging in a MeldableHeap]{Merging \ensuremath{\ensuremath{\ensuremath{\mathit{h_1}}}} and \ensuremath{\ensuremath{\ensuremath{\mathit{h_2}}}} is done by merging \ensuremath{\ensuremath{\ensuremath{\mathit{h_2}}}} with one of
  \ensuremath{\ensuremath{\ensuremath{\mathit{h_1}}.\ensuremath{\mathit{left}}}} or \ensuremath{\ensuremath{\ensuremath{\mathit{h_1}}.\ensuremath{\mathit{right}}}}.}
  \figlabel{meldable-merge}
\end{figure}
In the next section, we show that \ensuremath{\ensuremath{\mathrm{merge}(\ensuremath{\mathit{h_1}},\ensuremath{\mathit{h_2}})}} runs in $O(\log \ensuremath{\ensuremath{\ensuremath{\mathit{n}}}})$
expected time, where \ensuremath{\ensuremath{\ensuremath{\mathit{n}}}} is the total number of elements in \ensuremath{\ensuremath{\ensuremath{\mathit{h_1}}}} and \ensuremath{\ensuremath{\ensuremath{\mathit{h_2}}}}.
With access to a \ensuremath{\ensuremath{\mathrm{merge}(\ensuremath{\mathit{h_1}},\ensuremath{\mathit{h_2}})}} operation, the \ensuremath{\ensuremath{\mathrm{add}(\ensuremath{\mathit{x}})}} operation is easy.  We create a new node \ensuremath{\ensuremath{\ensuremath{\mathit{u}}}} containing \ensuremath{\ensuremath{\ensuremath{\mathit{x}}}} and then merge \ensuremath{\ensuremath{\ensuremath{\mathit{u}}}} with the root of our heap:
%importing \codeimport{ods/MeldableHeap.add(x)}: 
\begin{oframed}
\begin{flushleft}
\hspace*{1em} \ensuremath{\mathrm{add}(\ensuremath{\mathit{x}})}\\
\hspace*{1em} \hspace*{1em} \ensuremath{\ensuremath{\mathit{u}} \gets  \ensuremath{\mathrm{new\_node}(\ensuremath{\mathit{x}})}}\\
\hspace*{1em} \hspace*{1em} \ensuremath{\ensuremath{\mathit{r}} \gets  \ensuremath{\mathrm{merge}(\ensuremath{\mathit{u}}, \ensuremath{\mathit{r}})}}\\
\hspace*{1em} \hspace*{1em} \ensuremath{\ensuremath{\mathit{r}}.\ensuremath{\ensuremath{\mathit{parent}} \gets  nil}}\\
\hspace*{1em} \hspace*{1em} \ensuremath{\ensuremath{\mathit{n}} \gets  \ensuremath{\ensuremath{\mathit{n}} + 1}}\\
\hspace*{1em} \hspace*{1em} {\color{black} \textbf{return}} \ensuremath{\ensuremath{\mathit{true}}}\\
\end{flushleft}
\end{oframed}
This takes $O(\log (\ensuremath{\ensuremath{\ensuremath{\mathit{n}}}}+1)) = O(\log \ensuremath{\ensuremath{\ensuremath{\mathit{n}}}})$ expected time.
The \ensuremath{\ensuremath{\mathrm{remove}()}} operation is similarly easy.  The node we want to remove
is the root, so we just merge its two children and make the result the root:
%importing \codeimport{ods/MeldableHeap.remove()}: 
\begin{oframed}
\begin{flushleft}
\hspace*{1em} \ensuremath{\mathrm{remove}()}\\

\hspace*{1em} \hspace*{1em} \ensuremath{\ensuremath{\mathit{x}} \gets  \ensuremath{\ensuremath{\mathit{r}}.x}}\\
\hspace*{1em} \hspace*{1em} \ensuremath{\ensuremath{\mathit{r}} \gets  \ensuremath{\mathrm{merge}(\ensuremath{\mathit{r}}.\ensuremath{\mathit{left}}, \ensuremath{\mathit{r}}.\ensuremath{\mathit{right}})}}\\
\hspace*{1em} \hspace*{1em} {\color{black} \textbf{if}} \ensuremath{\ensuremath{\mathit{r}} \ne nil} {\color{black} \textbf{then}}  \ensuremath{\ensuremath{\mathit{r}}.\ensuremath{\ensuremath{\mathit{parent}} \gets  nil}}\\
\hspace*{1em} \hspace*{1em} \ensuremath{\ensuremath{\mathit{n}} \gets  \ensuremath{\ensuremath{\mathit{n}} - 1}}\\
\hspace*{1em} \hspace*{1em} {\color{black} \textbf{return}} \ensuremath{\ensuremath{\mathit{x}}}\\
\end{flushleft}
\end{oframed}
Again, this takes $O(\log \ensuremath{\ensuremath{\ensuremath{\mathit{n}}}})$ expected time.
Additionally, a MeldableHeap can implement many other operations in
$O(\log \ensuremath{\ensuremath{\ensuremath{\mathit{n}}}})$ expected time, including:
\begin{itemize}
\item \ensuremath{\ensuremath{\mathrm{remove}(\ensuremath{\mathit{u}})}}: remove the node \ensuremath{\ensuremath{\ensuremath{\mathit{u}}}} (and its key \ensuremath{\ensuremath{\ensuremath{\mathit{u}}.\ensuremath{\mathit{x}}}}) from the heap.
\item \ensuremath{\ensuremath{\mathrm{absorb}(\ensuremath{\mathit{h}})}}: add all the elements of the MeldableHeap \ensuremath{\ensuremath{\ensuremath{\mathit{h}}}} to this heap, emptying \ensuremath{\ensuremath{\ensuremath{\mathit{h}}}} in the process.
\end{itemize}
Each of these operations can be implemented using a constant number of
\ensuremath{\ensuremath{\mathrm{merge}(\ensuremath{\mathit{h_1}},\ensuremath{\mathit{h_2}})}} operations that each take $O(\log \ensuremath{\ensuremath{\ensuremath{\mathit{n}}}})$ expected time.
\subsection{Analysis of \ensuremath{\ensuremath{\mathrm{merge}(\ensuremath{\mathit{h_1}},\ensuremath{\mathit{h_2}})}}}
The analysis of \ensuremath{\ensuremath{\mathrm{merge}(\ensuremath{\mathit{h_1}},\ensuremath{\mathit{h_2}})}} is based on the analysis of a random walk
in a binary tree.  A \emph{random walk} in a binary tree starts at the
root of the tree.  At each step in the random walk, a coin is tossed and,
depending on the result of this coin toss, the walk proceeds to the left
or to the right child of the current node.  The walk ends when it falls
off the tree (the current node becomes \ensuremath{\ensuremath{\ensuremath{\mathit{nil}}}}).
The following lemma is somewhat remarkable because it does not depend
at all on the shape of the binary tree:
\begin{lem}\lemlabel{tree-random-walk}
The expected length of a random walk in a binary tree with \ensuremath{\ensuremath{\ensuremath{\mathit{n}}}} nodes is at most \ensuremath{\log \ensuremath{(\ensuremath{\mathit{n}}+1)}}.
\end{lem}
\begin{proof}
The proof is by induction on \ensuremath{\ensuremath{\ensuremath{\mathit{n}}}}. In the base case, $\ensuremath{\ensuremath{\ensuremath{\mathit{n}}}}=0$ and the walk
has length $0=\log (\ensuremath{\ensuremath{\ensuremath{\mathit{n}}}}+1)$.  Suppose now that the result is true for
all non-negative integers $\ensuremath{\ensuremath{\ensuremath{\mathit{n}}}}'< \ensuremath{\ensuremath{\ensuremath{\mathit{n}}}}$.
Let $\ensuremath{\ensuremath{\ensuremath{\mathit{n}}}}_1$ denote the size of the root's left subtree, so that
$\ensuremath{\ensuremath{\ensuremath{\mathit{n}}}}_2=\ensuremath{\ensuremath{\ensuremath{\mathit{n}}}}-\ensuremath{\ensuremath{\ensuremath{\mathit{n}}}}_1-1$ is the size of the root's right subtree.  Starting at
the root, the walk takes one step and then continues in a subtree of
size $\ensuremath{\ensuremath{\ensuremath{\mathit{n}}}}_1$ or $\ensuremath{\ensuremath{\ensuremath{\mathit{n}}}}_2$.  By our inductive hypothesis, the expected
length of the walk is then
\[
    \E[W] = 1 + \frac{1}{2}\log (\ensuremath{\ensuremath{\ensuremath{\mathit{n}}}}_1+1) + \frac{1}{2}\log (\ensuremath{\ensuremath{\ensuremath{\mathit{n}}}}_2+1)  \enspace , 
\] 
since each of $\ensuremath{\ensuremath{\ensuremath{\mathit{n}}}}_1$ and $\ensuremath{\ensuremath{\ensuremath{\mathit{n}}}}_2$ are less than $\ensuremath{\ensuremath{\ensuremath{\mathit{n}}}}$.  Since $\log$
is a concave function, $\E[W]$ is maximized when $\ensuremath{\ensuremath{\ensuremath{\mathit{n}}}}_1=\ensuremath{\ensuremath{\ensuremath{\mathit{n}}}}_2=(\ensuremath{\ensuremath{\ensuremath{\mathit{n}}}}-1)/2$.
%To maximize this,
%over all choices of $\ensuremath{\ensuremath{\ensuremath{\mathit{n}}}}_1\in[0,\ensuremath{\ensuremath{\ensuremath{\mathit{n}}}}-1]$, we take the derivative and obtain
%\[
%    (\E[W])' = \frac{1}{2}(c/\ensuremath{\ensuremath{\ensuremath{\mathit{n}}}}_1 - c/(\ensuremath{\ensuremath{\ensuremath{\mathit{n}}}}-\ensuremath{\ensuremath{\ensuremath{\mathit{n}}}}_1-1)) \enspace , 
%\]
%which is equal to 0 when $\ensuremath{\ensuremath{\ensuremath{\mathit{n}}}}_1 = (\ensuremath{\ensuremath{\ensuremath{\mathit{n}}}}-1)/2$.  We can establish that
%this is a maximum fairly easily, so
Therefore,
 the expected number of steps taken by the random walk is 
\begin{align*}
    \E[W] 
    & = 1 + \frac{1}{2}\log (\ensuremath{\ensuremath{\ensuremath{\mathit{n}}}}_1+1) + \frac{1}{2}\log (\ensuremath{\ensuremath{\ensuremath{\mathit{n}}}}_2+1) \\
   & \le  1 + \log ((\ensuremath{\ensuremath{\ensuremath{\mathit{n}}}}-1)/2+1) \\
   & =  1 + \log ((\ensuremath{\ensuremath{\ensuremath{\mathit{n}}}}+1)/2) \\
   & =  \log (\ensuremath{\ensuremath{\ensuremath{\mathit{n}}}}+1)  \enspace . \qedhere 
\end{align*}
\end{proof}
We make a quick digression to note that, for readers who know a little
about information theory, the proof of \lemref{tree-random-walk} can
be stated in terms of entropy.  
\begin{proof}[Information Theoretic Proof of \lemref{tree-random-walk}]
Let $d_i$ denote the depth of the $i$th external node and recall that a
binary tree with \ensuremath{\ensuremath{\ensuremath{\mathit{n}}}} nodes has \ensuremath{\ensuremath{\ensuremath{\mathit{n}}+1}} external nodes.  The probability
of the random walk reaching the $i$th external node is exactly
$p_i=1/2^{d_i}$, so the expected length of the random walk is given by
\[
   H=\sum_{i=0}^{\ensuremath{\ensuremath{\ensuremath{\mathit{n}}}}} p_id_i
    =\sum_{i=0}^{\ensuremath{\ensuremath{\ensuremath{\mathit{n}}}}} p_i\log\left(2^{d_i}\right)
    = \sum_{i=0}^{\ensuremath{\ensuremath{\ensuremath{\mathit{n}}}}}p_i\log({1/p_i})
\]
The right hand side of this equation is easily recognizable as the
entropy of a probability distribution over $\ensuremath{\ensuremath{\ensuremath{\mathit{n}}}}+1$ elements.  A basic
fact about the entropy of a distribution over $\ensuremath{\ensuremath{\ensuremath{\mathit{n}}}}+1$ elements is that
it does not exceed $\log(\ensuremath{\ensuremath{\ensuremath{\mathit{n}}}}+1)$, which proves the lemma.
\end{proof}
With this result on random walks, we can now easily prove that the
running time of the \ensuremath{\ensuremath{\mathrm{merge}(\ensuremath{\mathit{h_1}},\ensuremath{\mathit{h_2}})}} operation is $O(\log \ensuremath{\ensuremath{\ensuremath{\mathit{n}}}})$.
\begin{lem}
  If \ensuremath{\ensuremath{\ensuremath{\mathit{h_1}}}} and \ensuremath{\ensuremath{\ensuremath{\mathit{h_2}}}} are the roots of two heaps containing $\ensuremath{\ensuremath{\ensuremath{\mathit{n}}}}_1$
  and $\ensuremath{\ensuremath{\ensuremath{\mathit{n}}}}_2$ nodes, respectively, then the expected running time of
  \ensuremath{\ensuremath{\mathrm{merge}(\ensuremath{\mathit{h_1}},\ensuremath{\mathit{h_2}})}} is at most $O(\log \ensuremath{\ensuremath{\ensuremath{\mathit{n}}}})$, where $\ensuremath{\ensuremath{\ensuremath{\mathit{n}}}}=\ensuremath{\ensuremath{\ensuremath{\mathit{n}}}}_1+\ensuremath{\ensuremath{\ensuremath{\mathit{n}}}}_2$.
\end{lem}
\begin{proof}
  Each step of the merge algorithm takes one step of a random walk,
  either in the heap rooted at \ensuremath{\ensuremath{\ensuremath{\mathit{h_1}}}} or the heap rooted at \ensuremath{\ensuremath{\ensuremath{\mathit{h_2}}}}.
  %, depending on whether $\ensuremath{\ensuremath{\ensuremath{\mathit{h_1}}.\ensuremath{\mathit{x}}}} < \ensuremath{\ensuremath{\ensuremath{\mathit{h_2}}.\ensuremath{\mathit{x}}}}$ or not.
  The algorithm terminates when either of these two random walks fall
  out of its corresponding tree (when $\ensuremath{\ensuremath{\ensuremath{\mathit{h_1}}}}=\ensuremath{\ensuremath{\ensuremath{\mathit{nil}}}}$ or $\ensuremath{\ensuremath{\ensuremath{\mathit{h_2}}}}=\ensuremath{\ensuremath{\ensuremath{\mathit{nil}}}}$).
  Therefore, the expected number of steps performed by the merge algorithm
  is at most
  \[
     \log (\ensuremath{\ensuremath{\ensuremath{\mathit{n}}}}_1+1) + \log (\ensuremath{\ensuremath{\ensuremath{\mathit{n}}}}_2+1) \le 2\log \ensuremath{\ensuremath{\ensuremath{\mathit{n}}}} \enspace . \qedhere
  \]
\end{proof}
\subsection{Summary}
The following theorem summarizes the performance of a MeldableHeap:
\begin{thm}\thmlabel{meldableheap}
  A MeldableHeap implements the (priority) Queue interface.
  A MeldableHeap supports the operations \ensuremath{\ensuremath{\mathrm{add}(\ensuremath{\mathit{x}})}} and \ensuremath{\ensuremath{\mathrm{remove}()}}
  in $O(\log \ensuremath{\ensuremath{\ensuremath{\mathit{n}}}})$ expected time per operation.
\end{thm}
\section{Discussion and Exercises}
The implicit representation of a complete binary tree as an array,
or list, seems to have been first proposed by Eytzinger \cite{e1590}.
He used this representation in books containing pedigree family trees
\index{pedigree family tree}%
of noble families.  The BinaryHeap data structure described here was
first introduced by Williams \cite{w64}.
The randomized MeldableHeap data structure described here appears
to have first been proposed by Gambin and Malinowski \cite{gm98}.
Other meldable heap implementations exist, including 
leftist heaps \cite[Section~5.3.2]{c72,k97v3},
\index{leftist heap}%
\index{heap!leftist}%
binomial heaps \cite{v78},
\index{binomial heap}%
\index{heap!binomial}%
Fibonacci heaps \cite{ft87}, 
\index{Fibonacci heap}%
\index{heap!Fibonacci}%
pairing heaps \cite{fsst86},\
\index{pairing heap}%
\index{heap!pairing}%
 and skew heaps \cite{st83}, 
\index{skew heap}%
\index{heap!skew}%
although none of these are as simple as the MeldableHeap
structure.
Some of the above structures also support a \ensuremath{\ensuremath{\mathrm{decrease\_key}(\ensuremath{\mathit{u}},\ensuremath{\mathit{y}})}} operation
\index{decreaseKey@\ensuremath{\ensuremath{\mathrm{decrease\_key}(\ensuremath{\mathit{u}},\ensuremath{\mathit{y}})}}}%
in which the value stored at node \ensuremath{\ensuremath{\ensuremath{\mathit{u}}}} is decreased to \ensuremath{\ensuremath{\ensuremath{\mathit{y}}}}.  (It is a
pre-condition that $\ensuremath{\ensuremath{\ensuremath{\mathit{y}}}}\le\ensuremath{\ensuremath{\ensuremath{\mathit{u}}.\ensuremath{\mathit{x}}}}$.)  In most of the preceding structures,
this operation can be supported in $O(\log \ensuremath{\ensuremath{\ensuremath{\mathit{n}}}})$ time by removing node
\ensuremath{\ensuremath{\ensuremath{\mathit{u}}}} and adding  \ensuremath{\ensuremath{\ensuremath{\mathit{y}}}}.  However, some of these structures can implement
\ensuremath{\ensuremath{\mathrm{decrease\_key}(\ensuremath{\mathit{u}},\ensuremath{\mathit{y}})}} more efficiently.  In particular, \ensuremath{\ensuremath{\mathrm{decrease\_key}(\ensuremath{\mathit{u}},\ensuremath{\mathit{y}})}}
takes $O(1)$ amortized time in Fibonacci heaps and $O(\log\log \ensuremath{\ensuremath{\ensuremath{\mathit{n}}}})$
amortized time in a special version of pairing heaps \cite{e09}.
This more efficient \ensuremath{\ensuremath{\mathrm{decrease\_key}(\ensuremath{\mathit{u}},\ensuremath{\mathit{y}})}} operation has applications in
speeding up several graph algorithms, including Dijkstra's shortest path
algorithm \cite{ft87}.
\begin{exc}
  Illustrate the addition of the values 7 and then 3 to the BinaryHeap
  shown at the end of \figref{heap-insert}.
\end{exc}
\begin{exc}
  Illustrate the removal of the next two values (6 and 8) on the
  BinaryHeap shown at the end of \figref{heap-remove}.
\end{exc}
\begin{exc}
  Implement the \ensuremath{\ensuremath{\mathrm{remove}(\ensuremath{\mathit{i}})}} method, that removes the value stored in
  \ensuremath{\ensuremath{\ensuremath{\mathit{a}}[\ensuremath{\mathit{i}}]}} in a BinaryHeap.  This method should run in $O(\log \ensuremath{\ensuremath{\ensuremath{\mathit{n}}}})$ time.
  Next, explain why this method is not likely to be useful.
\end{exc}
\begin{exc}\exclabel{general-eytzinger}
  \index{tree!$d$-ary}%
  A $d$-ary tree is a generalization of a binary tree in which each
  internal node has $d$ children.  Using Eytzinger's method it is also
  possible to represent complete $d$-ary trees using arrays.  Work out
  the equations that, given an index \ensuremath{\ensuremath{\ensuremath{\mathit{i}}}}, determine the index of \ensuremath{\ensuremath{\ensuremath{\mathit{i}}}}'s
  parent and each of \ensuremath{\ensuremath{\ensuremath{\mathit{i}}}}'s $d$ children in this representation.
\end{exc}
\begin{exc}
  \index{DaryHeap@DaryHeap}%
  Using what you learned in \excref{general-eytzinger}, design and
  implement a \emph{DaryHeap}, the $d$-ary generalization of a
  BinaryHeap. Analyze the running times of operations on a DaryHeap
  and test the performance of your DaryHeap implementation against
  that of the BinaryHeap implementation given here.
\end{exc}
\begin{exc}
  Illustrate the addition of the values 17 and then 82 in the
  MeldableHeap \ensuremath{\ensuremath{\ensuremath{\mathit{h_1}}}} shown in \figref{meldable-merge}.  Use a coin to
  simulate a random bit when needed.
\end{exc}
\begin{exc}
  Illustrate the removal of the next two values (4 and 8) in the
  MeldableHeap \ensuremath{\ensuremath{\ensuremath{\mathit{h_1}}}} shown in \figref{meldable-merge}.  Use a coin to
  simulate a random bit when needed.
\end{exc}
\begin{exc}
  Implement the \ensuremath{\ensuremath{\mathrm{remove}(\ensuremath{\mathit{u}})}} method, that removes the node \ensuremath{\ensuremath{\ensuremath{\mathit{u}}}} from
  a MeldableHeap.  This method should run in $O(\log \ensuremath{\ensuremath{\ensuremath{\mathit{n}}}})$ expected time.
\end{exc}
\begin{exc}
  Show how to find the second smallest value in a BinaryHeap or
  MeldableHeap in constant time.
\end{exc}
\begin{exc}
  Show how to find the $k$th smallest value in a BinaryHeap or
  MeldableHeap in $O(k\log k)$ time.  (Hint: Using another heap
  might help.)
\end{exc}
\begin{exc}
  Suppose you are given \ensuremath{\ensuremath{\ensuremath{\mathit{k}}}} sorted lists, of total length \ensuremath{\ensuremath{\ensuremath{\mathit{n}}}}.  Using
  a heap, show how to merge these into a single sorted list in $O(n\log
  k)$ time.  (Hint: Starting with the case $k=2$ can be instructive.)
\end{exc}
